\SourceLeaf{373}%
und schließlich auch denselben Krümmungskreis, wenn zusätzlich gilt
\[
\frac{rqq-2spq+tpp}{r'q'q'-2s'p'q'+t'p'p'}=\frac{p^3}{p'^3}=\frac{q^3}{q'^3}.
\]
Daraus lässt sich der Krümmungskreis leicht bestimmen, wenn man annimmt, dass eine der beiden Kurven ein Kreis ist.

Sind also $x,y$ orthogonale Koordinaten, so ist der Mittelpunkt des Krümmungskreises durch die Gleichungen bestimmt
\[
\begin{aligned}
x&=\xi+\frac{p(pp+qq)}{rqq-2spq+tpp}\\[5pt]
y&=\eta-\frac{q(pp+qq)}{rqq-2spq+tpp}.
\end{aligned}
\]
\needspace{7\baselineskip}
\begin{center}BEMERKUNG.\end{center}

Die vorstehende Notiz findet sich auf dem Umschlag des Werks \textit{Elementorum analyseos infinitorum Joanni Andreae Segneri Pars I.} Halle 1761 und stammt höchstwahrscheinlich aus dem Ende des achtzehnten Jahrhunderts.\par\nobreak\hfill\textsc{Stäckel.}\par
